%!TEX program = lualatex
% SPDX-License-Identifier: CC-BY-SA-4.0
\documentclass[11pt]{article}

\usepackage{iftex}
\ifLuaTeX\else
  \PackageError{vintage-latex}{This example requires LuaLaTeX}
    {Run lualatex, not pdflatex or xelatex.}
\fi

\usepackage{fontspec}
\usepackage{xcolor}
\usepackage{geometry}
\usepackage[tracking=true,protrusion=false,expansion=false]{microtype}

\defaultfontfeatures{Renderer=HarfBuzz}
\setmainfont{EB Garamond}[Numbers={OldStyle,Proportional}]
\definecolor{cataloguepaper}{HTML}{F5EBD8}
\definecolor{printerink}{HTML}{292119}
\pagecolor{cataloguepaper}
\color{printerink}
\geometry{a5paper,top=19mm,bottom=21mm,left=20mm,right=20mm}
\pagestyle{empty}
\setlength{\parindent}{0pt}
\setlength{\emergencystretch}{.7em}

\newcommand{\entry}[4]{%
  \par\medskip
  \noindent\makebox[2.25em][l]{\textbf{#1}}%
  \textit{#2}\dotfill\textls[35]{\textsc{#3}}\par
  \hangindent=2.25em\hangafter=0
  \small #4\par
}

\begin{document}

\begin{center}
  {\small\textls[120]{\textsc{Museum of Practical Natural History}}}\par
  \vspace{.35em}
  {\LARGE Catalogue of Minerals}\par
  \vspace{.2em}
  {\itshape Cabinet B: metallic and vitreous specimens}\par
  \vspace{.55em}
  \rule{\linewidth}{1pt}\par
  \vspace{.16em}
  \rule{\linewidth}{.3pt}
\end{center}

\vspace{.8em}

\entry{B.1}{Galena}{Cumberland}{Lead-grey cubes with three perfect cleavages;
the fresh fracture is bright, but the exposed face has dulled. Specific gravity
recorded as 7.48.}

\entry{B.2}{Chalcopyrite}{Cornwall}{Brass-yellow massive specimen with an
iridescent tarnish. The streak is greenish black and should be observed on an
unglazed tile.}

\entry{B.3}{Hematite}{Elba}{Botryoidal surface over a compact red-brown mass.
The powdered streak distinguishes it from magnetite even where the exterior is
nearly black.}

\entry{B.4}{Fluorite}{Derbyshire}{Pale violet interpenetrating cubes upon
limestone. Four cleavage directions yield octahedral fragments under a light
blow.}

\entry{B.5}{Quartz}{St. Gotthard}{A clear six-sided prism terminated by paired
rhombohedra. Minute chlorite inclusions mark an earlier stage of growth.}

\vspace{1em}
\begin{center}
  \rule{.62\linewidth}{.35pt}\par
  \vspace{.45em}
  {\footnotesize\textls[70]{\textsc{Key to the entries}}}\par
  \vspace{.35em}
  \begin{tabular}{@{}rl@{}}
    Roman type & observed description\\
    \textit{Italic type} & accepted specimen name\\
    \textsc{Small capitals} & locality or collection\\
    Old-style figures & measurements and catalogue numbers
  \end{tabular}
\end{center}

\vfill
\begin{center}
  {\scriptsize Revised after inspection of the cabinet, 1906}
\end{center}

\end{document}